\documentclass[11pt]{article}
\usepackage{amsmath}
\usepackage{amsfonts}
\usepackage{amsthm}
\usepackage[utf8]{inputenc}
\usepackage[margin=0.75in]{geometry}

\title{CSC111 Winter 2026 Project 1}
\author{Jack Anderson, Efren Medina Arias, Diego Jose Figarella Urbaneja}
\date{\today}

\begin{document}
\maketitle

\section*{Running the game}
Run the game by executing \verb|adventure.py|. The game uses the \verb|tkinter| module to open a file dialog when loading game files, though this requires no additional installation. It should be noted that in \verb|commands.py|, there is a Python TA flag for the unused import of \verb|Location| and \verb|Event|, though these classes \textit{are} used in doctests in the file.

\section*{Game Map}
Game map is provided below:

\begin{verbatim}
 -1    -1    18    -1    -1    -1    41    -1
 -1    -1    17    -1    -1    37    40    -1
 -1    -1    16    36    35    34    38    39
 -1    -1    15    -1    -1    33    -1    32
 14    13    12    30    29    27    28    31
 -1    -1    11    -1    -1    26    -1    -1
 -1    -1    10    25    24    23    -1    -1
 -1    -1     9    -1    -1    20    21    22
 -1    -1     8    -1    -1    19    -1    -1
 -1    -1     7     6     5     2     3     4
 -1    -1    -1    -1    -1     1    -1    -1
 -1    -1    -1    -1    -1     0    -1    -1
\end{verbatim}
North is up. There are two tunnels: 22$\leftrightarrow$31 and 32$\leftrightarrow$33.

Starting location is: 1, which is outside of the dorm, and the target location for all items is 0, which is the dorm itself.

\section*{Game solution}
The player must collect three items (USB Drive from location 7, Laptop Charger from location 10, Lucky Mug from location 4) and bring them all back to location 0 (the dorm room). In order to get back to location 0, the player \textbf{must} solve the puzzle at Hart House to unlock teleportation, otherwise the player will run out of steps. One efficient winning path is:

\begin{verbatim}
go north, go west, go west, enter wallberg, take usb drive, exit wallberg north,
go north, enter myhal north, take laptop charger, go north, go north, go west,
enter mclennan west, exit mclennan east, go east, go east, go east, go east,
go east, enter gerstein east, enter tunnel to hart house, go north,
code 137202111, teleport 4, take lucky mug,
teleport 0
\end{verbatim}

\section*{Lose condition(s)}
Description of how to lose the game: The player has a limited number of steps (25 by default). Each movement command decrements the step counter by 1. If the player runs out of steps before collecting all three items and returning to the dorm (location 0), the game ends with the message ``You've run out of time! The 1pm deadline has passed.''

List of commands (example that results in a loss --- wander until steps run out):
\begin{verbatim}
go north, go south, go north, go south, go north, go south, go north, go south,
go north, go south, go north, go south, go north, go south, go north, go south,
go north, go south, go north, go south, go north, go south, go north, go south,
go north
\end{verbatim}

Which parts of your code are involved in this functionality:
\begin{itemize}
    \item \verb|event_logger.py|: \verb|GameState.use_step()|, \verb|GameState.has_steps()| --- these manage the step counter and check if the player has steps remaining.
    \item \verb|commands.py|: \verb|CommandHandler.move_to_location()| --- calls \verb|use_step()| on each movement.
    \item \verb|adventure.py|: the main game loop checks \verb|game.state.has_steps()| each iteration and terminates the game if no steps remain.
\end{itemize}

\section*{Inventory}

\begin{enumerate}
\item All location IDs that involve items in the game: 7 (Wallberg Building), 10 (Myhal), 4 (Coffee Shop), and 0 (Campus One Dorm Room --- the target location for all items).

\item Item data:
\begin{enumerate}
    \item For Item 1:
    \begin{itemize}
    \item Item name: USB Drive
    \item Item start location ID: 7 (Wallberg Building)
    \item Item target location ID: 0 (Campus One Dorm Room)
    \end{itemize}
    \item For Item 2:
    \begin{itemize}
    \item Item name: Laptop Charger
    \item Item start location ID: 10 (Myhal)
    \item Item target location ID: 0 (Campus One Dorm Room)
    \end{itemize}
    \item For Item 3:
    \begin{itemize}
    \item Item name: Lucky Mug
    \item Item start location ID: 4 (Coffee Shop)
    \item Item target location ID: 0 (Campus One Dorm Room)
    \end{itemize}
\end{enumerate}

    \item Example commands to pick up and drop an item:
\begin{verbatim}
['go north', 'go west', 'go west', 'enter wallberg', 'take usb drive',
 'inventory', 'drop usb drive', 'inventory']
\end{verbatim}
    This navigates to Wallberg (location 7), picks up the USB Drive with \verb|take usb drive|, checks inventory, drops it with \verb|drop usb drive|, and checks inventory again.

    \item Which parts of your code are involved in handling the \verb|inventory| command:
    \begin{itemize}
        \item \verb|game_entities.py|: \verb|Item| dataclass --- represents each item with name, description, start/current/target position.
        \item \verb|event_logger.py|: \verb|GameState.take_item()|, \verb|GameState.drop_item()| --- move items between the ground list and the inventory list.
        \item \verb|commands.py|: \verb|CommandHandler.take()|, \verb|CommandHandler.drop()|, \verb|CommandHandler.inventory_menu()| --- parse user input, find matching items, and call the appropriate \verb|GameState| methods or display the inventory.
    \end{itemize}
\end{enumerate}

\section*{Score}
\begin{enumerate}

    \item Players earn score in the following ways:
    \begin{itemize}
        \item \textbf{+5 points} for visiting a new (unvisited) location for the first time. The first location where the player earns score is location 1 (Kings Circle), which is awarded automatically at game start. After that, moving to location 2 with \verb|go north| earns another 5 points.
        \item \textbf{+10 points} for picking up an item (via \verb|take|).
        \item \textbf{+100 points} for solving the puzzle at Hart House (location 39) by entering the correct combination.
    \end{itemize}
    First score increase: the player scores 5 points at the starting location 1 upon beginning the game (since it is the first unvisited location).

    Commands to demonstrate score increase: \verb|['go north', 'score']|. After moving north (visiting location 2, +5 points), the player can type \verb|score| to see their total.

    \item List assigned to \verb|scores_demo|:
\begin{verbatim}
scores_demo = ['go north', 'score']
\end{verbatim}

    \item Which parts of your code are involved in handling the \verb|score| functionality:
    \begin{itemize}
        \item \verb|event_logger.py|: \verb|GameState| class stores the \verb|score| attribute; \verb|GameState.take_item()| adds 10 to the score when an item is picked up.
        \item \verb|commands.py|: \verb|CommandHandler.get_score()| displays the current score to the player.
        \item \verb|adventure.py|: the main game loop awards +5 score each time a new (unvisited) location is entered; \verb|AdventureGame.check_puzzle_combination()| awards +100 for solving the puzzle.
    \end{itemize}
\end{enumerate}

\section*{Enhancements}

\begin{enumerate}

    \item GameState class and state tracking with EventList
    \begin{itemize}
       \item \textbf{Description:} Implemented a new class, \verb|GameState|, to store all relevant game state information that will change over the course of the game. This included implementing functionality for deep-copying, through implementing methods on \verb|Item|, \verb|Location| and \verb|GameState| itself. This paved the way for implementing the undo command later on.
       \item \textbf{Complexity level:} Medium
       \item \textbf{Reasons:} This feature warrants the medium complexity level because it required a complete redesign of the entire game's codebase, and a significant amount of planning. Planning stages included considering which attributes of \verb|AdventureGame| were mutable and important to state, and implementing these into \verb|GameState|. Then, we needed to augment \verb|EventList| to store instances of the game state for the fututre \verb|undo| command. Following an initial implementation, we had to begin a long process of debugging and studying mutability, aliases and references accross all parts of the codebase which accessed the game state in some capacity. This resulted in implementing methods for deep-copying the game state, which cascaded into implementing copy methods on all game entities in \verb|game_entities.py|, in addition to \verb|GameState|.
       \item \textbf{Parts of code involved:}
       \begin{itemize}
        \item \verb|game_entities.py|: \verb|Item.copy| and \verb|Location.copy| methods
        \item \verb|event_logger.py|: \verb|GameState| class and associated methods, including \verb|GameState.copy()|. \verb|EventList| class was augmented to store instances of \verb|GameState|.
        \item \verb|adventure.py|: augmented \verb|AdventureGame| to include a state attribute, to which we assign an instance of \verb|GameState|.
       \end{itemize}
       \item \textbf{Demo commands:} Any command demonstrates this, as all actions update or reference the player's current state.
    \end{itemize}

    \item Undo Function
    \begin{itemize}
        \item \textbf{Description:} The undo command allows the player to delete the last event in \verb|EventList| and restore the \verb|GameState| to the previous event state. When the player types \verb|undo|, the game restores a deep copy of the \verb|GameState| from the previous event in the \verb|EventList| linked list, and removes the most recent event node. This effectively rewinds all the \verb|GameState| attributes like score, locations visited, steps, etc. If there are no moves to undo, a message is displayed instead.

        \item \textbf{Complexity level:} Medium

        \item \textbf{Reasons:} Implementing undo required the creation of a deep copy of the entire game state at every event. This meant implementing \verb|copy()| methods on the \verb|Item|, \verb|Location|, and \verb|GameState| classes so that each event in the \verb|EventList| stores an independent snapshot of the game. Without deep copying, since \verb|GameState| is a constantly mutated object, then due to aliasing, the game wouldn't be able to go back to the saved ``copy'' of \verb|GameState|, which is just a reference to the current \verb|GameState|. The undo logic itself leverages the doubly-linked list structure of \verb|EventList| since we are able to simply access the previous node of the last node of \verb|EventList|. We also had to ensure proper synchronization between the \verb|GameState| in \verb|CommandHandler| and \verb|AdventureGame|. We did a reference reassignment in the main loop after an undo. Additionally, undo is classified as a ``menu command'' so that it does not itself create a new event node in the log.

        \item \textbf{Parts of the code involved:}
        \begin{itemize}
            \item \verb|game_entities.py|: \verb|Location.copy()|, \verb|Item.copy()| --- produces independent copies of game entities, so snapshots are immutable.
            \item \verb|event_logger.py|: \verb|GameState.copy()| --- deep copies the entire game state (inventory, locations dict, items list, score, steps, etc.). \verb|EventList.remove_last_event()| --- removes the tail node of the doubly-linked list and resets the new tail's \verb|next_command| and \verb|next| pointers.
            \item \verb|adventure.py|: the main loop re-syncs \verb|game.state = command_handler.game_state| after each command, ensuring the undo state propagates correctly.
        \end{itemize}

        \item \textbf{Demo commands:}
\begin{verbatim}
enhancement1_demo = ['go north', 'go east', 'undo', 'undo']
\end{verbatim}
    \end{itemize}

    \item Game Saving and Loading
    \begin{itemize}
        \item \textbf{Description:} The player can save the game at any point by typing \verb|save|. This exports a JSON file that saves all the data in \verb|EventList|. The player can use the JSON file to load the previously saved game on startup by choosing ``load'' at the entry menu. The load dialogue uses the \verb|tkinter| library so that the user can easily upload the JSON file to load the game.

        \item \textbf{Complexity level:} High

        \item \textbf{Reasons:} This portion of the project was highly complicated to implement, in addition to the \verb|GameState| function and associated functionality. It required the creation of methods to serialize the complex datatypes (\verb|Item|, \verb|Location|, \verb|Event|, \verb|EventList|) into more primitive data types (list and dict), in order to be written to a JSON file. In addition, we needed to implement the reverse: functions that read serialized JSON and rebuilt \verb|Event|s, \verb|Item|s and \verb|Location|s to be injected into the \verb|AdventureGame| class at the beginning of the game.

        \item \textbf{Parts of the code involved:}
        \begin{itemize}
            \item \verb|event_logger.py|: \verb|GameState.to_dictionary()| --- serializes the game state to a JSON-ready dict. \verb|EventList.to_list()| --- serializes the entire event linked list to a list of dicts.
            \item \verb|commands.py|: \verb|CommandHandler.save_game()| --- copies the current state, converts it to a dictionary, appends the serialized event log, and writes the result to \verb|savegame.json|.
            \item \verb|adventure.py|: \verb|rebuild(events, event_list, adventure_game)| --- deserializes a list of event dicts back into \verb|Event| objects with fully reconstructed \verb|GameState|s and reconnects them in the \verb|EventList|. \verb|load_game()| --- opens a \verb|tkinter| file dialog, reads the JSON save file, calls \verb|rebuild()|, and restores \verb|adventure_game.state| to the last event's state. \verb|game_entry()| --- presents the player with the load/new menu at startup.
        \end{itemize}

        \item \textbf{Demo commands:}
\begin{verbatim}
enhancement2_demo = ['go north', 'go west', 'go west', 'enter wallberg',
                     'take usb drive', 'save']
\end{verbatim}
    \end{itemize}

    \item Command Handler System
    \begin{itemize}
        \item \textbf{Description:} Implemented a command handler class and associated methods to streamline iteration and creation of future commands. The \verb|CommandHandler| class is instantiated once per game, and takes an instance of \verb|EventList| and \verb|GameState|. Given the considerations above for implementation of \verb|GameState|, and the addition of \verb|EventList| for activity history, the command handler effectively has access to the entire mutable game through references.

        \item \textbf{Complexity level:} Medium

        \item \textbf{Reasons:} This portion of the project was moderately complicated to implement, as we had to consider ways to avoid circularity in our code, and remain organized at the same time. We had to consider the role that references and mutability played in assigning \verb|AdventureGame|'s \verb|GameState| to \verb|CommandHandler.game_state|, especially for implementing commands within \verb|CommandHandler| which dealt with previous states, such as the undo command. Additionally, we implemented the \verb|_parse_command()| method, which separated commands and arguments, and a \verb|get_allowed()| function to dynamically generate a list of all allowed commands in a certain location. The small amount of complexity in this implementation came from having to work with \verb|Location.available_commands| as strings including spaces (e.g.\ \verb|'go east'|), which are important as whole strings, in the same context of commands where the root and arguments have distinct meanings. In addition, the command handler's ability to parse plaintext commands and produce outputs in the game was especially useful for implementing the simulation very quickly.

        \item \textbf{Parts of the code involved:}
        \begin{itemize}
            \item \verb|commands.py|: contains \verb|CommandHandler| class with \verb|_parse_command()|, \verb|run_command()|, and individual methods for running each command.
            \item \verb|adventure.py|: inside \verb|__main__| we instantiate \verb|CommandHandler| with the relevant game objects.
        \end{itemize}

        \item \textbf{Demo commands:} Any command demonstrates this, as they're all routed through the handler.
    \end{itemize}

    \item Puzzle and Code
    \begin{itemize}
        \item \textbf{Description:} Implemented a puzzle feature where, once at Hart House, the player has to input a lock combination (137202111) in order to enter a room and meet Bert, the ghost of Hart House, who gives the player the ability to teleport to any location in the game. The puzzle provides a hint of what the numbers could be or where one could find additional hints (it references MAT137, McLennan's MP202, and CSC111). Players can interact with the puzzle using the code command and successfully entering 'code 137202111' changes the \texttt{unlimited\_transportations} bool of the current \texttt{GameState} to \texttt{True}, adds 100 points to your score and 10 steps to your total steps left.
    
        \item \textbf{Complexity level:} Medium
    
        \item \textbf{Reasons:} This feature was moderately complex to implement, mostly because of its contextual nature. The puzzle appears at Hart House and 'code' is available as a command only if the puzzle hasn't unlocked teleportation yet, so we had to constantly check the player's location state and use dynamic command generation by calling the method \verb|CommandHandler.get_allowed()|. This section also required constant parsing of the numeric values from the code command in order to accurately update the step limit and score.
    
        \item \textbf{Parts of the code involved:}
        \begin{itemize}
            \item \verb|commands.py|: contains \verb|CommandHandler.code()|, which is the method that parses the number arguments, validates the input format, calls the puzzle check, and display success or failure messages. It also includes the \verb|get_allowed()| method that appends code to all available commands when in Hart House.
            \item \verb|adventure.py|: Inside \verb|__main__|, our game loop displays the puzzle hint at location 39.
            \item \verb|event_logger.py|: contains \verb|GameState.check_puzzle_combination()| which validates the code parsed from the player command and updates the game state.
        \end{itemize}
    
        \item \textbf{Demo commands:} Same commands as game solution demonstrate puzzle and code.
    \end{itemize}
    
    \item Teleportation
        \begin{itemize}
            \item \textbf{Description:} Implemented a teleportation command that allows the player to move to any in game location using one step. It is only accessible once the player has solved the Hart House puzzle. Teleport can be invoked with a numeric argument (i.e., the ID of the location the player wants to teleport to), and it can be invoked without any argument, which then prompts the game to show a list of all possible location IDs. Teleportation is needed to win the game since the initial amount of steps given is not enough to reach all 3 items and return to the dorm.

            \item \textbf{Complexity level:} Low
        
            \item \textbf{Reasons:} teleportation was a relatively straightforward  feature once we had already done the dynamic command generation system. Our implementation checks whether \verb|unlimited_transporations| is True and simply updates the player's current location ID with the selected ID. We only required basic conditional statements and extending the available commands list using the \verb|get_allowed()| function we had already designed.
        
            \item \textbf{Parts of the code involved:}
            \begin{itemize}
                \item \verb|commands.py|: contains \verb|CommandHandler.teleport_anywhere()|, which either lists all locations (when no arguments are passed) or validates input, calling \verb|move_to_location()|. Also, \verb|get_allowed()| method that appends teleport to all available commands when \\ \texttt{unlimited\_transportations} is True.
                \item \verb|event_logger.py|: contains \verb|unlimited_transportations| bool instance attribute in GameState, representing if teleport has been unlocked yet.
            \end{itemize}
        
            \item \textbf{Demo commands:} Same commands as game solution demonstrate the teleportation enhancement.
        \end{itemize}
\end{enumerate}


\end{document}
